\chapter{Detecting classes and parameter-space arguments}\label{sec:sweepout-topology}
Let $M^{n+1}$ be smooth, closed and connected, $n\geq1$. We use the mod-two boundary component, the class $\bar\lambda\in H^1(\mathcal Z_n^0;\mathbb Z_2)$, and admissibility of Definition \ref{def:p-width}. Almgren's identification of the detecting class remains external. The arguments below use ordinary singular cohomology with $\mathbb Z_2$ coefficients and its cellular computation for finite complexes. Sources are \cite{LMN}, Sections 2.5--2.6, and \cite{MNIndex}, Section 5.

\begin{lemma}[Pulling back a detecting family]\label{lem:sweepout-pullback}
Let $\Phi:X\to\mathcal Z_n^0$ be an admissible family on a finite cubical complex, and let $f:Y\to X$ be continuous with $Y$ another finite cubical complex. Then $\Phi\circ f$ is flat-continuous, has no concentration, and has maximal mass at most that of $\Phi$. It is a $p$-sweepout exactly when
\[
 f^*(\Phi^*\bar\lambda^p)\ne0.
\]
\end{lemma}
\begin{proof}
Continuity follows by composition. Every slice of the composite is a slice of the original family, so its maximal mass and every uniform small-ball supremum are bounded by those of $\Phi$. Naturality of cohomological pullback gives $(\Phi\circ f)^*(\bar\lambda^p)=f^*(\Phi^*\bar\lambda^p)$. This proves the exact detection criterion; the geometric bounds alone do not imply it.
\end{proof}

\begin{proposition}[A parameter-dimension obstruction]\label{prop:sweepout-parameter-dimension}
A $p$-sweepout on a finite cubical complex $X$ requires $\dim X\geq p$. Having dimension at least $p$ is not sufficient for a family to be a $p$-sweepout.
\end{proposition}
\begin{proof}
The cellular cochain groups above $\dim X$ vanish, and cellular cohomology agrees with singular cohomology. Hence $H^p(X;\mathbb Z_2)=0$ when $p>\dim X$, contradicting detection. Conversely a constant family has zero pullback in every positive cohomological degree, regardless of the dimension of its parameter space. The cellular comparison theorem is the external algebraic-topology input.
\end{proof}

\begin{proposition}[The p-skeleton retains a p-sweepout]\label{prop:sweepout-skeleton}
If $\Phi:X\to\mathcal Z_n^0$ is an admissible $p$-sweepout on a finite cubical complex, its restriction to the $p$-skeleton $X^{(p)}$ is an admissible $p$-sweepout with no larger maximal mass.
\end{proposition}
\begin{proof}
Let $i:X^{(p)}\hookrightarrow X$ be inclusion. The relative cellular cochain complex for $(X,X^{(p)})$ has no cochains in degrees at most $p$, so $H^p(X,X^{(p)};\mathbb Z_2)=0$. The long exact sequence of the pair therefore makes $i^*:H^p(X;\mathbb Z_2)\to H^p(X^{(p)};\mathbb Z_2)$ injective. Thus the nonzero class $\Phi^*\bar\lambda^p$ remains nonzero on the skeleton. Lemma \ref{lem:sweepout-pullback} gives the analytic admissibility and mass bound. This reduces the parameter dimension of a detecting family without making a claim about the index of a limiting hypersurface.
\end{proof}

\begin{proposition}[Detection is invariant under flat homotopy]\label{prop:sweepout-homotopy-detection}
Suppose $\Phi_0,\Phi_1:X\to\mathcal Z_n^0$ are flat-homotopic and both independently satisfy the bounded-mass and no-concentration conditions. Then $\Phi_0$ is a $p$-sweepout if and only if $\Phi_1$ is. The homotopy itself need not satisfy a uniform no-concentration estimate for this cohomological conclusion. Compare \cite{LMN}, Section 2.6.
\end{proposition}
\begin{proof}
A homotopy into the cycle space with its flat topology induces the same pullback on singular cohomology at its two endpoints. Thus $\Phi_0^*\bar\lambda^p=\Phi_1^*\bar\lambda^p$. Nonvanishing is equivalent, and the separately assumed analytic conditions make each endpoint admissible. This argument cannot be used to infer no concentration for an endpoint from flat homotopy alone; Example \ref{ex:flat-continuous-concentrating-family} explains the missing control.
\end{proof}

\begin{proposition}[A detecting fixed class bounds the spectral infimum]\label{prop:fixed-class-bounds-spectrum}
Let $X$ have a finite cubical structure and let $\Phi:X\to\mathcal Z_n^0$ be an $\mathbf F$-continuous $p$-sweepout. Let $\Pi$ be its flat-homotopy class among $\mathbf F$-continuous maps on $X$, as in Definition \ref{def:homotopy-width}, using a triangulation when expressing $X$ as a simplicial complex. Then
\[
 \omega_p(M,g)\leq\mathbf L(\Pi).
\]
There is no assertion of equality for an arbitrary such class.
\end{proposition}
\begin{proof}
Every $\Psi\in\Pi$ has bounded maximal mass and no concentration by Proposition \ref{prop:f-continuity-no-concentration}. It detects the same cup power as $\Phi$ by Proposition \ref{prop:sweepout-homotopy-detection}. Thus every competitor in the infimum defining $\mathbf L(\Pi)$ is also a competitor for $\omega_p$, whose infimum allows other domains and classes as well. This inclusion of competitor sets gives the inequality. Triangulation changes neither the underlying space nor its cohomology or maps; no interpolation between different cycle topologies is used in this comparison.
\end{proof}

\begin{example}[The same slices can lose detection after reparametrization]\label{ex:double-traversal-sweepout}
Let $\Phi:S^1\to\mathcal Z_n^0$ be an admissible one-sweepout, with $S^1$ given a finite cubical structure. The family $\Psi(z)=\Phi(z^2)$ has the same image, the same maximal mass and the same small-ball suprema, but is not a one-sweepout.
\end{example}
\begin{proof}
The degree-two map $f(z)=z^2$ is surjective, so the sets of slices are identical. On $H^1(S^1;\mathbb Z_2)$ its pullback is multiplication by two, hence zero. Therefore $\Psi^*\bar\lambda=f^*(\Phi^*\bar\lambda)=0$. The elementary degree computation can be seen by traversing the oriented circle twice and reducing the resulting homology class modulo two. It shows why covering the same geometric slices does not preserve the topological obstruction.
\end{proof}

\begin{lemma}[Cup-product vanishing on an open cover]\label{lem:cup-product-open-cover}
Let $X$ be a finite complex, let $U_1,\ldots,U_N$ be open subsets covering $X$, and let $\alpha_i\in H^{d_i}(X;\mathbb Z_2)$ with $d_i\geq1$. If $\alpha_i|_{U_i}=0$ for every $i$, then $\alpha_1\smile\cdots\smile\alpha_N=0$. This is the topological step behind the localization argument described in \cite{LMN}, pp.~935--937.
\end{lemma}
\begin{proof}
The long exact sequence of $(X,U_i)$ lifts $\alpha_i$ to a class $\widetilde\alpha_i\in H^{d_i}(X,U_i;\mathbb Z_2)$. The relative cup product for open, hence excisive, pairs places the product of these lifts in
\[
 H^{d_1+\cdots+d_N}(X,U_1\cup\cdots\cup U_N;\mathbb Z_2).
\]
Its image in absolute cohomology is the desired product by naturality. Since the union is $X$, the relative group is zero. The long exact sequence and relative cup product are the explicit external algebraic-topology inputs; this argument supplies the deduction from them.
\end{proof}

\begin{proposition}[A common parameter above several cost thresholds]\label{prop:cup-length-cost-thresholds}
Let $\beta\in H^1(X;\mathbb Z_2)$ on a finite complex satisfy $\beta^p\ne0$. Suppose $p_i\geq1$ and $\sum_{i=1}^N p_i\leq p$. Let $c_i:X\to[0,\infty)$ be continuous, let $a_i>0$, and assume $\beta^{p_i}$ restricts to zero on $U_i=\{c_i<a_i\}$. Then some $x\in X$ satisfies $c_i(x)\geq a_i$ for every $i$. If $F(x)\geq\sum_i c_i(x)$ for all $x$, then $\sup_XF\geq\sum_i a_i$.
\end{proposition}
\begin{proof}
If the $U_i$ covered $X$, Lemma \ref{lem:cup-product-open-cover} would give $\beta^{\sum_i p_i}=0$. Multiplying by the remaining cup power, or using this equality directly if the sum is $p$, would give $\beta^p=0$, a contradiction. Thus a point lies outside every $U_i$ and satisfies all thresholds. The bound for $F$ follows at that point. Applying this abstract lemma to cycle masses in subdomains requires additional work: those costs need not be continuous in the flat topology, and restriction need not give finite-mass relative boundaries. These are precisely the issues recorded in \cite{LMN}, p.~937; their geometric resolution is not proved by the present topological lemma.
\end{proof}
